\documentclass[11pt,a4paper]{article}
\usepackage[margin=2.5cm]{geometry}
\usepackage{amsmath,amssymb,amsthm}
\usepackage{booktabs}
\usepackage{enumitem}
\usepackage{hyperref}
\usepackage[T1]{fontenc}
\usepackage[utf8]{inputenc}

\title{Piecewise Birth--Death Model (PWBD):\\Implementation and Recovery Experiments}
\author{PhyloBayesSynth Project}
\date{June 2026}

\begin{document}
\maketitle

%======================================================================
\section{Grammar Extension: PWBD}
%======================================================================

We extend the context-free grammar with a recursive \textbf{Piecewise Birth--Death} (PWBD) combinator.
The full grammar for model expressions ($N_3$) is:
\begin{align*}
  N_3 &\;::=\; (\texttt{CRB}\; P)
         \;\mid\; (\texttt{CRBD}\; P\; P)
         \;\mid\; (\texttt{TDB}\; P\; P)
         \;\mid\; (\texttt{TDBD}\; P\; P\; P) \\
      &\;\;\;\mid\; (\texttt{PWBD}\; P\; N_3\; N_3)
         \;\mid\; (\texttt{BAMM}\; P\; N_3\; N_3)
\end{align*}
The PWBD node takes one parameter $\tau_{\mathrm{frac}} \in (0,1)$ and two child model expressions.
The parameter $\tau_{\mathrm{frac}}$ specifies the breakpoint as a \emph{relative} fraction of the current time interval:
given an interval $[\tau_{\mathrm{lo}},\, \tau_{\mathrm{hi}}]$, the breakpoint is
$\tau_b = \tau_{\mathrm{lo}} + \tau_{\mathrm{frac}} \cdot (\tau_{\mathrm{hi}} - \tau_{\mathrm{lo}})$.
The first child governs $[\tau_{\mathrm{lo}},\, \tau_b]$ (near-present) and the second governs $[\tau_b,\, \tau_{\mathrm{hi}}]$ (near-root).

Because PWBD is recursive, nesting produces arbitrarily many segments with heterogeneous model types:
\[
  \underbrace{(\texttt{PWBD}\; 0.4\; (\texttt{CRB}\; 0.5)\;
    (\texttt{PWBD}\; 0.6\; (\texttt{CRBD}\; 0.3\; 0.1)\; (\texttt{TDB}\; 0.8\; {-}0.1)))}_{\text{3 segments: CRB $\mid$ CRBD $\mid$ TDB}}
\]
Maximum nesting depth is controlled by the configuration parameter \texttt{max\_pwbd\_nesting}~(default~3).
BAMM is retained in the grammar but its prior probability is set to zero by default.

%======================================================================
\section{Closed-Form Likelihood via $\Lambda$-Parameterization}
%======================================================================

For any model where the ratio $\varepsilon = \mu/\lambda$ is constant within a segment
(covering CRB, CRBD, TDB, and TDBD), the Riccati extinction ODE
\[
  \frac{dE}{d\tau} = \lambda(\tau)\,(E-1)(E-\varepsilon), \quad E(\tau_0)=E_0
\]
is separable and admits an analytical solution parameterized by the cumulative birth rate
$\Lambda(\tau_0,\tau) = \int_{\tau_0}^{\tau}\lambda(s)\,ds$.
Define $\varphi = (1-\varepsilon)\Lambda$ and
\[
  V(\Lambda) = (E_0-\varepsilon) - (E_0-1)\,e^{\varphi}.
\]
Then the extinction probability, survival probability, and $D$-propagation factor are:
\begin{align}
  E &= \frac{(E_0-\varepsilon) - \varepsilon(E_0-1)\,e^{\varphi}}{V(\Lambda)}, \label{eq:E}\\[4pt]
  \log(1-E) &= \log(1{-}\varepsilon) + \log(1{-}E_0) + \varphi - \log|V|, \label{eq:logsurvival}\\[4pt]
  \log \frac{D(\Lambda_p)}{D(\Lambda_c)}
    &= (1{-}\varepsilon)(\Lambda_p - \Lambda_c) + 2\log|V(\Lambda_c)| - 2\log|V(\Lambda_p)|. \label{eq:Dfactor}
\end{align}
The PWBD likelihood is computed by:
(i)~flattening the nested PWBD tree into an ordered list of $K$ intervals,
(ii)~forward-sweeping $E$ across intervals (chaining initial conditions at breakpoints via \eqref{eq:E}),
(iii)~splitting each branch at breakpoints and computing $D$-factors per segment via \eqref{eq:Dfactor},
and (iv)~conditioning on root survival via \eqref{eq:logsurvival}.
The total cost is $O(n \cdot K)$---fully closed-form with no ODE integration.

%======================================================================
\section{Data-Driven PWBD Proposal}
%======================================================================

A key challenge is that PWBD expressions sampled from the prior have random parameters
far from the likelihood peak, causing near-zero acceptance in MH.
We address this with a \textbf{data-driven proposal} mechanism:

\begin{enumerate}[nosep]
  \item \textbf{Breakpoint scan}: evaluate 8 candidate $\tau_{\mathrm{frac}}$ values uniformly in $[0.15,\, 0.85]$.
  \item \textbf{Per-segment rate estimation}: for each candidate,
        split branching events and branches into two groups at the breakpoint.
        Estimate $\hat\lambda_k = n_k / S_k$ (events per total branch length)
        and test $\mu$-candidates at fractions $\{0, 0.1, 0.3, 0.5\} \times \hat\lambda_k$.
  \item \textbf{Selection}: choose the $(\tau_{\mathrm{frac}},\, \lambda_1, \mu_1,\, \lambda_2, \mu_2)$
        combination with highest likelihood.
  \item \textbf{Jitter}: perturb all parameters with log-normal noise ($\sigma=0.15$) to maintain exploration.
\end{enumerate}
Additionally, after any structural (non-local) proposal, we run 30 steps of greedy hill-climbing on parameters.

The MH move probabilities are: local parameter perturbation (65\%),
data-driven PWBD proposal (15\%), and structural regeneration from the prior (20\%).

%======================================================================
\section{BIC Complexity Penalty}
%======================================================================

Without explicit complexity control, nested PWBD can overfit by adding segments that each capture noise.
We add a BIC-style penalty to the MH acceptance ratio.
For a structural move from expression $E$ with $k$ free parameters to $E'$ with $k'$ parameters,
the log-acceptance ratio becomes:
\[
  \log\alpha = \underbrace{(\ell' - \ell)}_{\text{likelihood ratio}}
             + \underbrace{\log\frac{|A_E|}{|A_{E'}|}}_{\text{size correction}}
             + \underbrace{\log\frac{\pi(E')}{\pi(E)}}_{\text{prior ratio}}
             - \underbrace{w_{\mathrm{BIC}}\,(k'-k)\,\log n}_{\text{complexity penalty}}
\]
where $n$ is the number of tips and $w_{\mathrm{BIC}}$ is a tunable weight (default $0.5$, corresponding to the standard BIC).
This ensures that each additional parameter must improve the log-likelihood by at least $\frac{1}{2}\log n$ nats to be accepted.

%======================================================================
\section{Recovery Experiments}
%======================================================================

We test the framework with two experiments using cumulative-posterior SMC
(60~particles, 25 rejuvenation moves per round, 8 trees of 100~tips each).

%----------------------------------------------------------------------
\subsection{Experiment 1: CRBD Control}
%----------------------------------------------------------------------

\textbf{True model}: $(\texttt{CRBD}\; 0.3\; 0.05)$.

\smallskip\noindent
\begin{tabular}{lrl}
\toprule
\textbf{Metric} & \textbf{Value} & \textbf{Expression} \\
\midrule
True model cum-logL     & 1122.07  & $(\texttt{CRBD}\; 0.300\; 0.050)$ \\
Best posterior cum-logL  & 1125.30  & $(\texttt{CRBD}\; 0.346\; 0.097)$ \\
\midrule
\multicolumn{3}{l}{\textbf{Posterior model distribution (60 particles):}} \\
\quad CRB  & 41/60 & (68.3\%) \\
\quad CRBD & 19/60 & (31.7\%) \\
\quad PWBD &  0/60 & (0.0\%) \\
\bottomrule
\end{tabular}

\smallskip
The posterior is concentrated on CRB and CRBD.
No particles adopt PWBD, confirming that the BIC penalty successfully prevents overfitting
when the true model has no rate shift.
The recovered parameters $(\hat\lambda=0.346,\; \hat\mu=0.097)$ are close to the true values.

%----------------------------------------------------------------------
\subsection{Experiment 2: PWBD Recovery}
%----------------------------------------------------------------------

\textbf{True model}: $(\texttt{PWBD}\; 0.5\; (\texttt{CRB}\; 0.3)\; (\texttt{CRBD}\; 0.8\; 0.15))$
--- pure birth near the present, birth--death with extinction near the root.

\smallskip\noindent
\begin{tabular}{lrl}
\toprule
\textbf{Metric} & \textbf{Value} & \textbf{Expression} \\
\midrule
True model cum-logL     & 1727.42  & $(\texttt{PWBD}\; 0.5\; (\texttt{CRB}\; 0.3)\; (\texttt{CRBD}\; 0.8\; 0.15))$ \\
Best posterior cum-logL  & 1832.47  & $(\texttt{PWBD}\; 0.614\; (\texttt{CRBD}\; 0.765\; 0.156)\;$ \\
                         &          & \quad $(\texttt{CRBD}\; 1.613\; 1.582))$ \\
\midrule
\multicolumn{3}{l}{\textbf{Posterior model distribution (60 particles):}} \\
\quad PWBD & 60/60 & (100.0\%) \\
\bottomrule
\end{tabular}

\smallskip
All 60 particles converge to PWBD, correctly identifying the piecewise structure.
The recovered near-present segment has $\hat\mu_1/\hat\lambda_1 = 0.20$
(low extinction), while the near-root segment has $\hat\mu_2/\hat\lambda_2 = 0.98$
(near-critical, heavy extinction)---qualitatively matching the true pattern of
zero extinction near the present and significant extinction near the root.
The breakpoint $\hat\tau_{\mathrm{frac}} = 0.614$ is reasonably close to the true value of 0.5.

Note that the recovered \emph{absolute} parameter values differ from the truth
($\hat\lambda_1 = 0.765$ vs.\ true $\lambda_1 = 0.3$),
reflecting the well-known non-identifiability of absolute rates from reconstructed phylogenies:
the MLE parameters can differ from the generating parameters while producing higher likelihood on any finite realization.
The qualitative structure (breakpoint location, extinction contrast) is the identifiable and correctly recovered signal.

\end{document}
